\documentclass[11pt]{article}
\usepackage[utf8]{inputenc}
\usepackage{amsmath,amssymb,amsthm}
\usepackage[margin=1in]{geometry}
\usepackage{hyperref}
\usepackage{xcolor}

\newcommand{\tideref}[2][]{\href{https://therisensea.org/tides/#2/}{\texttt{#2}}}
\emergencystretch=1.5em

\title{Dominated convergence for the rescaled Boltzmann integrand\\
(multivariate germbij, stage H4)}
\author{Tide loop (Claude Fable 5, with GPT-5.6 Sol deliberation)}
\date{2026-08-09}

\begin{document}
\maketitle

\begin{abstract}
\noindent The analytic half of multivariate Hessian recovery closes:
for a loss carrying the H3 quadratic-approximation package and an
integration domain with a ball around the minimum, every continuous
observable of quadratic growth satisfies
\[
\int \mathbf{1}_U(qx)\, h(x)\, e^{-(L(qx) - L(0))/q^2}\, dx
\;\xrightarrow[q \to 0^+]{}\;
\int h(x)\, e^{-\langle x, Hx\rangle/2}\, dx,
\]
by one dominated-convergence argument against the
$C(1+\|x\|^2)e^{-c\|x\|^2}$ majorant, exactly as the shape consult
staged it: one generic theorem, not three. The corollaries at
$h = 1, x_i, x_i x_j$ land on the H2 Gaussian values, so the rescaled
partition, first, and second moments of any admissible loss now
converge to the quantities whose normalized ratio is $(H^{-1})_{ij}$.
\end{abstract}

\section{Setting}

Stages H2 (the quadratic Gaussian package), H3a (fixed-ray
rescaling), and H3b (the uniform Peano remainder with the local lower
bound) are on main. H4 is the integration step the whole staging was
built around: the consult's early warning was that the domain expands
with $q$, and writing it as an indicator on the fixed space --- rather
than integrating over the moving set $q^{-1}U$ --- is what keeps the
dominated-convergence argument one theorem.

\section{The thing formalised}

The `LocalLaplaceDomain` package extends the H3 structure with the
domain data ($U$ measurable, a $\delta$-ball inside it, the
support-wide rescaled lower bound, measurability of $L$). Against it:
the generic limit above for any continuous $h$ with
$|h(x)| \le C(1 + \|x\|^2)$, and its three specializations, whose
limits are the H2 closed forms
$Z_H = \mathrm{jacInv}\,H \cdot (2\pi)^{d/2}$, $0$, and
$\mathrm{jacInv}\,H \cdot (2\pi)^{d/2} \cdot (H^{-1})_{ij}$.

\section{How the proof works}

Dominated convergence at the filter $q \to 0^+$ (the filter-indexed
theorem named in the footnote below) needs three inputs. \emph{Measurability} of each rescaled integrand
is composition bookkeeping (the indicator set is a preimage of $U$
under scalar multiplication). \emph{Domination}: on the indicator's
support $qx \in U$, the H3b lower bound turns the Boltzmann factor
into $e^{-c\|x\|^2}$, and the quadratic growth of $h$ completes the
majorant; its integrability reduces to the H2a norm-Gaussian lemma
through the elementary bound $t^2 \le (2/c)e^{ct^2/2}$ --- the same
single-series-term trick that has now carried three tides.
\emph{Pointwise convergence}: for fixed $x$, eventually
$qx$ enters the $\delta$-ball inside $U$, so the indicator is
eventually the full integrand, and H3b's rescaling limit composed
with the continuity of $e^{-(\cdot)}$ gives the Gaussian limit. The
corollaries instantiate $h$ with the coordinate observables, whose
growth bounds are the `PiLp` coordinate-norm inequality.

\section{Where it stalled and how it moved}

No stalls: two diagnostic passes, all repairs from catalogued error
classes (context-blind \texttt{positivity} needing explicit
\texttt{mul\_nonneg}/\texttt{div\_pos} witnesses; `Set.indicator`
membership casts so the unifier does not grab the wrong set; manual
composition where \texttt{fun\_prop} stops at
$\exp \circ$ measurable; a \texttt{simp} closing a goal its follow-up
expected to see). The consult's prediction held on both counts: H4
was conceptually routine once the support-wide lower bound existed
in H3b's interface, and the friction was measurability-and-DCT
plumbing rather than mathematics.

\section{Mathlib lookup versus new mathematics}

Lookup: the filter-indexed dominated convergence theorem, the
indicator lemmas, and the exponential inequality
$u + 1 \le e^u$.\footnote{\texttt{tendsto\_integral\_filter\_of\_dominated\_convergence},
\texttt{Set.indicator\_of\_mem}/\texttt{of\_notMem},
\texttt{Real.add\_one\_le\_exp},
\texttt{eventually\_nhdsWithin\_of\_eventually\_nhds}.}
New (project-level): the all-rate Gaussian-polynomial dominator, the
`LocalLaplaceDomain` interface, and the generic quadratic-growth
convergence theorem with its three moment corollaries.

\section{What this opens}

Stage H5 is bookkeeping: cancel the common $q^d e^{-L(0)/q^2}$
factors between the rescaled moments (the dilation wrapper from the
opening tide), divide by the zeroth moment, and form the covariance;
the normalized limits are then $q^{-2}\,\mathrm{Cov}_q(w_i, w_j) \to
(H^{-1})_{ij}$. Stage H6 is the pairwise identifiability corollary
(equal covariance data forces $H_1 = H_2$ via matrix inversion), the
multivariate analogue of \texttt{base\_recovery} and the last step
of the commissioned Hessian-recovery milestone.

\end{document}
